\documentclass[a4paper,12pt]{article}
\usepackage[T2A]{fontenc}
\usepackage[utf8]{inputenc}
\usepackage[russian]{babel}
\usepackage{amsmath,amssymb,mathtools}
\usepackage{hyperref}
\usepackage{enumitem}
\usepackage{booktabs}
\usepackage{geometry}
\geometry{top=2cm, bottom=2cm, left=2.5cm, right=2.5cm}

\title{Конспект лекции 7: \\ Языковые модели (LLM)}
\author{Липкин С.М.}
\date{}

\begin{document}
\maketitle

\section{Постановка задачи языкового моделирования}

Языковое моделирование --- фундаментальная задача NLP. Модель учится предсказывать распределение вероятностей последовательностей токенов.

\subsection{Формальное определение}

Для последовательности токенов $w_1, w_2, \dots, w_n$ совместная вероятность раскладывается через цепное правило:

\[
P(w_1, w_2, \dots, w_n) = \prod_{t=1}^n P(w_t \mid w_{1}, \dots, w_{t-1}) = \prod_{t=1}^n P(w_t \mid w_{<t}).
\]

Задача языковой модели --- научиться аппроксимировать $P(w_t \mid w_{<t})$ для любого контекста.

\subsection{История}

\begin{itemize}
    \item \textbf{N-граммы:} считаем частоты $n$ последовательных токенов. Не работают на длинных контекстах.
    \item \textbf{RNN/LSTM:} обрабатывают последовательность шаг за шагом, но медленно и с затухающими градиентами.
    \item \textbf{Transformer (2017):} параллельная обработка всей последовательности через self-attention.
    \item \textbf{LLM (2018--настоящее время):} Transformer-архитектуры на масштабе миллиардов параметров.
\end{itemize}

\section{Encoder vs Decoder: два подхода}

Все современные LLM используют архитектуру Transformer, но по-разному: encoder-only для понимания, decoder-only для генерации, encoder-decoder для трансформаций.

\subsection{Encoder-only (BERT)}
\begin{itemize}
    \item Каждый токен видит полный контекст (бидирекционально).
    \item Нет каузальной маски --- attention на все токены.
    \item Используется: классификация, NER, QA, NLI.
    \item Примеры: BERT, RoBERTa, ALBERT, DistilBERT.
\end{itemize}

\subsection{Decoder-only (GPT)}
\begin{itemize}
    \item Каждый токен видит только предыдущие (causal mask).
    \item Авторегрессивная генерация: предсказываем следующий токен.
    \item Используется: генерация текста, диалог, код.
    \item Примеры: GPT, LLaMA, Mistral, Qwen.
\end{itemize}

\subsection{Encoder-Decoder (T5, BART)}
\begin{itemize}
    \item Encoder читает вход (бидирекционально).
    \item Decoder генерирует выход (авторегрессивно, с cross-attention к encoder).
    \item Используется: перевод, суммаризация, преобразование текста.
    \item Примеры: T5, BART, PEGASUS.
\end{itemize}

\section{Encoder-based модели: BERT}

BERT (Bidirectional Encoder Representations from Transformers) предложен Google в 2018 году (Devlin et al.).

\subsection{Архитектура}
BERT использует многослойный Transformer encoder. Стандартные конфигурации:
\begin{itemize}
    \item BERT-Base: 12 слоёв, 768 hidden, 12 heads = 110M параметров.
    \item BERT-Large: 24 слоя, 1024 hidden, 16 heads = 340M параметров.
\end{itemize}

\subsection{Предобучение}
BERT обучается на двух задачах одновременно:

\textbf{1. Masked Language Modeling (MLM):}
\begin{itemize}
    \item 15\% токенов выбираются для маскировки.
    \item Из выбранных: 80\% заменяются на [MASK], 10\% --- на случайный токен, 10\% --- остаются без изменений.
    \item Модель учится восстанавливать исходные токены.
    \item Это позволяет учитывать контекст слева и справа (бидирекциональность).
\end{itemize}

\textbf{2. Next Sentence Prediction (NSP):}
\begin{itemize}
    \item На вход подаются два предложения: A и B.
    \item Модель должна предсказать, является ли B следующим за A.
    \item [CLS]-токен аккумулирует информацию для классификации.
    \item Позже выяснилось, что NSP мало эффективен (RoBERTa его удалила).
\end{itemize}

\subsection{Fine-tuning}
После предобучения BERT донастраивается под конкретную задачу:

\begin{itemize}
    \item \textbf{Классификация:} [CLS] $\to$ линейный слой $\to$ softmax.
    \item \textbf{Token-level (NER, POS):} каждый выходной токен $\to$ линейный слой.
    \item \textbf{QA (SQuAD):} предсказание start/end позиции ответа.
\end{itemize}

\subsection{Варианты и развитие}
\begin{itemize}
    \item \textbf{RoBERTa:} больше данных, нет NSP, динамическая маскировка.
    \item \textbf{ALBERT:} разделение параметров (factorized embedding + cross-layer sharing).
    \item \textbf{DistilBERT:} дистилляция (40\% меньше, 97\% качества).
    \item \textbf{ELECTRA:} детекция заменённых токенов вместо MLM.
\end{itemize}

\section{Авторегрессионные модели: GPT}

GPT (Generative Pre-trained Transformer) предложен OpenAI в 2018 году.

\subsection{Архитектура и обучение}

GPT использует decoder-only Transformer с каузальным маскированием. Функция потерь --- стандартная кросс-энтропия для next-token prediction:

\[
\mathcal{L} = -\frac{1}{T} \sum_{t=1}^T \log P(x_t \mid x_{<t}; \theta).
\]

Обучение ведётся на огромных корпусах текстов (Common Crawl, WebText, Wikipedia, книги). Токенизация --- BPE (Byte Pair Encoding, Sennrich et al., 2016), которая разбивает слова на подслова, решая проблему OOV.

\subsection{Эволюция GPT}

\begin{itemize}
    \item \textbf{GPT-1 (2018):} 117M параметров, 12 слоёв. Предобучение + fine-tuning.
    \item \textbf{GPT-2 (2019):} 1.5B параметров, 48 слоёв. Zero-shot способности на многих задачах.
    \item \textbf{GPT-3 (2020):} 175B параметров, 96 слоёв. In-context learning (few-shot без fine-tuning).
    \item \textbf{GPT-4 (2023):} $\approx$1.7T параметров (предположительно MoE). Мультимодальность.
\end{itemize}

\subsection{Инференс}

Авторегрессивный инференс --- последовательная генерация:
\begin{enumerate}
    \item Подаём prompt.
    \item Получаем распределение $P(x_{t+1} \mid x_{1:t})$.
    \item Выбираем следующий токен (greedy, top-k, top-p, temperature).
    \item Добавляем к последовательности, повторяем.
\end{enumerate}

\textbf{KV Cache:} ключевая оптимизация. На каждом шаге пересчёт $K$ и $V$ для предыдущих токенов --- дорого. KV-кеш сохраняет $K$ и $V$ из предыдущих шагов, на новом шаге считаем только для нового токена. Снижает сложность с $O(n^2)$ до $O(n)$ на шаг.

\textbf{Стратегии декодирования:}
\begin{itemize}
    \item \textbf{Greedy:} $x_{t+1} = \arg\max P(\cdot \mid x_{1:t})$. Просто, но может зацикливаться.
    \item \textbf{Top-k:} ресемплинг из $k$ наиболее вероятных токенов.
    \item \textbf{Nucleus (top-p):} выбираем минимальное множество токенов с кумулятивной вероятностью $\ge p$.
    \item \textbf{Temperature:} $P_i \propto \exp(\text{logit}_i / T)$. $T \to 0$ --- argmax, $T \to \infty$ --- равномерное.
\end{itemize}

\subsection{Emergent abilities и Scaling Laws}

\textbf{Emergent abilities:} способности, которые не проявляются в маленьких моделях, но внезапно появляются при масштабировании. Примеры: in-context learning, chain-of-thought, арифметика, перевод.

\textbf{Scaling Laws (Kaplan et al., 2020):} loss модели следует степенному закону от числа параметров $N$ и объёма данных $D$:
\[
\text{Loss}(N, D) \propto N^{-\alpha} \cdot D^{-\beta},
\]
где $\alpha \approx 0.076$, $\beta \approx 0.103$. Вывод: увеличивать $N$ выгоднее, чем $D$.

\textbf{Chinchilla (Hoffmann et al., 2022):} DeepMind показала, что предыдущие результаты недооценивали важность данных. Оптимальное соотношение:
\[
N_{\text{opt}} \approx 2.1 \times 10^{9} \cdot C^{0.5},
\]
где $C$ --- compute budget в FLOPs. Для 175B модели нужно 3.7T токенов (у GPT-3 было 300B --- сильный under-training). Chinchilla (70B на 1.4T токенов) превзошла GPT-3 при меньшем размере.

\section{От GPT-3 к ChatGPT}

\subsection{Проблема raw GPT-3}
\begin{itemize}
    \item Умеет продолжать текст, но не следует инструкциям.
    \item Пример: \texttt{Переведи на русский: Hello} $\to$ \texttt{«Hello» может быть переведено...}
    \item Нужно: \texttt{Здравствуйте}
\end{itemize}

\subsection{Instruction Tuning (SFT)}
Собираем пары (инструкция, идеальный ответ) и дообучаем модель:
\begin{itemize}
    \item \textbf{FLAN (Wei et al., 2021):} 60+ NLP-задач, переформулированных как инструкции.
    \item \textbf{Alpaca (Stanford, 2023):} 52K инструкций от GPT-3.5 через Self-Instruct.
    \item \textbf{OpenAssistant:} краудсорсинговые датасеты диалогов.
    \item \textbf{Формат:} system / user / assistant роли в промпте.
\end{itemize}

\subsection{RLHF (Reinforcement Learning from Human Feedback)}
Трёхшаговый процесс от InstructGPT (Ouyang et al., 2022):

\textbf{Шаг 1: SFT} --- дообучение на инструкциях (как выше).

\textbf{Шаг 2: Reward Model (RM)} --- обучаем модель ранжировать ответы:
\begin{itemize}
    \item Генерируем несколько ответов на промпт.
    \item Человек ранжирует их (A $>$ B $>$ C).
    \item RM учится предсказывать предпочтения человека.
    \item Loss: Bradley-Terry model $P(A > B) = \sigma(r_A - r_B)$.
\end{itemize}

\textbf{Шаг 3: PPO (Proximal Policy Optimization)} --- оптимизируем LM:
\begin{itemize}
    \item Генерируем ответ $\to$ RM даёт награду $\to$ PPO обновляет LM.
    \item KL penalty: $\beta \cdot D_{KL}(\pi_{\text{PPO}} \parallel \pi_{\text{SFT}})$ --- чтобы не отойти далеко от SFT-модели.
    \item Итоговая loss: $\mathcal{L} = \mathbb{E}[-r(x, y)] + \beta \cdot D_{KL}(\pi_{\text{PPO}} \parallel \pi_{\text{SFT}})$.
\end{itemize}

\textbf{Ключевой результат InstructGPT:} 1.3B модель с RLHF превзошла 175B GPT-3 по оценкам людей. Alignment важнее размера.

\subsection{ChatGPT (ноябрь 2022)}
ChatGPT добавил к RLHF:
\begin{itemize}
    \item Диалоговый формат с multi-turn историей.
    \item Safety fine-tuning: отказ от вредных запросов.
    \item Контроль стиля и ограничений.
    \item Продвинутые методы инференса (latency/throughput оптимизации).
\end{itemize}

\section{Альтернативные подходы к alignment}

\begin{itemize}
    \item \textbf{DPO (Direct Preference Optimization, 2023):} явно задаём предпочтения без RM и RL. Оптимизируем напрямую: $\mathcal{L} = -\mathbb{E}[\log \sigma(\beta(\log \pi_\theta(y_w|x) - \log \pi_\theta(y_l|x)) - \beta \log(\pi_{\text{ref}}(y_w|x) / \pi_{\text{ref}}(y_l|x)))]$.
    \item \textbf{GRPO (Group Relative PPO):} DeepSeek. Без Critic-сети. Награда нормируется по группе ответов.
    \item \textbf{Constitutional AI:} модель сама генерирует и оценивает свои ответы по конституции (набору правил).
    \item \textbf{RLAIF:} Reward Model заменяется на LLM, дающую AI feedback.
\end{itemize}

\section{DeepSeek и Reasoning}

\subsection{DeepSeek-R1 (январь 2025)}
DeepSeek-R1 показала, что рассуждение (reasoning) можно обучить через RL без SFT:

\begin{enumerate}
    \item \textbf{Cold-start:} небольшой SFT на цепочках рассуждений.
    \item \textbf{GRPO (Group Relative Policy Optimization):} награда за правильный ответ + формат reasoning.
    \item \textbf{Self-evolution:} модель естественно учится более длинным и сложным рассуждениям.
    \item \textbf{Aha-moment:} модель «перепроверяет» свои рассуждения --- поведение, которого не было в обучении.
\end{enumerate}

\subsection{DeepSeek-V3 (декабрь 2024)}
\begin{itemize}
    \item 671B параметров, 37B активных (MoE с 256 экспертами).
    \item Обучен на 14.8T токенов. FP8 mixed-precision.
    \item Стоимость обучения: \$5.5M (в 10 раз дешевле аналогов).
    \item Использует Multi-Token Prediction (MTP) --- предсказание нескольких токенов сразу.
    \item Multi-head Latent Attention (MLA) --- сжатие KV кеша для эффективного инференса.
\end{itemize}

\section{Mixture of Experts (MoE)}

\subsection{Идея}
MoE заменяет плотные FFN-слои на спарсные наборы экспертов. Для каждого токена активируется только подмножество экспертов (обычно top-2):

\[
\text{MoE}(x) = \sum_{i=1}^E G(x)_i \cdot E_i(x), \quad G(x) = \text{softmax}(W_g x),
\]
где $G(x)$ --- gating network, $E_i(x)$ --- i-й эксперт (обычно FFN).

\subsection{Преимущества и недостатки}

\textbf{Преимущества:}
\begin{itemize}
    \item Больше параметров, те же FLOPs на токен.
    \item Качество близко к dense-модели с полным размером.
    \item Естественная специализация экспертов (разные эксперты для разных тем/задач).
\end{itemize}

\textbf{Недостатки:}
\begin{itemize}
    \item Все эксперты в памяти --- высокий VRAM.
    \item Load balancing: неравномерное распределение токенов по экспертам.
    \item Fine-tuning сложнее (overfitting, не все эксперты учатся).
    \item Коммуникационные overhead при распределённом обучении.
\end{itemize}

\subsection{Практические реализации}

\begin{itemize}
    \item \textbf{Mixtral 8x7B (Mistral, 2023):} 8 экспертов, top-2 активация. 47B всего / 12.9B активных.
    \item \textbf{DeepSeekMoE:} мелкозернистые эксперты (больше маленьких), shared expert для всех токенов, auxiliary loss-free load balancing через динамический bias.
    \item \textbf{Switch Transformer (Google, 2021):} top-1 активация (один эксперт на токен). Максимальная спарсность.
    \item \textbf{GShard:} распределённый MoE для машинного перевода.
\end{itemize}

\subsection{Технические детали DeepSeekMoE}
\begin{itemize}
    \item \textbf{Мелкозернистые эксперты:} 256 экспертов вместо 8-16, каждый меньше. Более гибкое распределение знаний.
    \item \textbf{Shared expert:} один общий FFN для всех токенов + routed эксперты для специализации.
    \item \textbf{Dynamic bias вместо aux loss:} регуляция загрузки экспертов через bias-терм в gating, без дополнительных потерь.
\end{itemize}

\section{Выводы}
\begin{itemize}
    \item Языковое моделирование --- предсказание следующего токена --- основа всех современных LLM.
    \item Encoder-only (BERT) для понимания, decoder-only (GPT) для генерации.
    \item Scaling Laws: модель и данные должны масштабироваться согласованно.
    \item RLHF выравнивает модель под человеческие предпочтения.
    \item DeepSeek показала, что reasoning можно получать через RL.
    \item MoE --- новый стандарт для масштабирования (больше параметров, те же FLOPs).
\end{itemize}

\end{document}
